\documentclass[11pt,a4paper]{article}
% Wspólny szablon dla kart supermocy (jedna karta = jedna strona A4)
% Użycie w karcie: \documentclass[11pt,a4paper]{article} % Wspólny szablon dla kart supermocy (jedna karta = jedna strona A4)
% Użycie w karcie: \documentclass[11pt,a4paper]{article} % Wspólny szablon dla kart supermocy (jedna karta = jedna strona A4)
% Użycie w karcie: \documentclass[11pt,a4paper]{article} \input{../preambula} \begin{document}

% Silnik tectonic (XeTeX): polskie znaki przez fontspec, nie przez inputenc/T1
\usepackage{fontspec}
\setmainfont{Arial}
\usepackage[polish]{babel}
\usepackage[a4paper,margin=1.5cm]{geometry}
\usepackage{amsmath,amssymb}
\usepackage{tikz}
\usetikzlibrary{arrows.meta,positioning,calc,shapes.geometric,decorations.pathreplacing}
\usepackage[most]{tcolorbox}
\usepackage{enumitem}
\usepackage{booktabs}
\usepackage{xcolor}

\pagestyle{empty}
\setlength{\parindent}{0pt}
\setlength{\parskip}{0.4em}

\definecolor{cel}{HTML}{1F4E79}
\definecolor{regula}{HTML}{E8F1FA}
\definecolor{pulapka}{HTML}{FDECEA}
\definecolor{okres}{HTML}{C0392B}
\definecolor{skonczony}{HTML}{1E8449}

% Ramka na regułę
\newtcolorbox{regula}[1][]{
  colback=regula, colframe=cel, boxrule=0.8pt, arc=2pt,
  left=6pt, right=6pt, top=4pt, bottom=4pt, title=#1,
  fonttitle=\bfseries, coltitle=white
}

% Ramka na pułapkę z kartkówki
\newtcolorbox{pulapka}[1][]{
  colback=pulapka, colframe=okres, boxrule=0.8pt, arc=2pt,
  left=6pt, right=6pt, top=4pt, bottom=4pt, title=#1,
  fonttitle=\bfseries, coltitle=white
}

% Nagłówek karty: numer, tytuł, cel
\newcommand{\naglowekKarty}[3]{%
  {\LARGE\bfseries\color{cel} Karta #1: #2}\\[2pt]
  {\large Cel: #3}\par
  \vspace{0.3em}\noindent\rule{\linewidth}{1pt}\par
}

% Pole na obliczenia ucznia: ramka z delikatną kratką, argument = wysokość
\newcommand{\miejsce}[1]{%
  \par\vspace{0.2em}\noindent
  \begin{tikzpicture}
    \draw[step=0.5cm, gray!25, very thin] (0,0) grid (\textwidth,#1);
    \draw[gray!60, rounded corners=2pt] (0,0) rectangle (\textwidth,#1);
    \node[anchor=north west, font=\scriptsize, text=gray] at (0.05,#1) {obliczenia};
  \end{tikzpicture}\par}
 \begin{document}

% Silnik tectonic (XeTeX): polskie znaki przez fontspec, nie przez inputenc/T1
\usepackage{fontspec}
\setmainfont{Arial}
\usepackage[polish]{babel}
\usepackage[a4paper,margin=1.5cm]{geometry}
\usepackage{amsmath,amssymb}
\usepackage{tikz}
\usetikzlibrary{arrows.meta,positioning,calc,shapes.geometric,decorations.pathreplacing}
\usepackage[most]{tcolorbox}
\usepackage{enumitem}
\usepackage{booktabs}
\usepackage{xcolor}

\pagestyle{empty}
\setlength{\parindent}{0pt}
\setlength{\parskip}{0.4em}

\definecolor{cel}{HTML}{1F4E79}
\definecolor{regula}{HTML}{E8F1FA}
\definecolor{pulapka}{HTML}{FDECEA}
\definecolor{okres}{HTML}{C0392B}
\definecolor{skonczony}{HTML}{1E8449}

% Ramka na regułę
\newtcolorbox{regula}[1][]{
  colback=regula, colframe=cel, boxrule=0.8pt, arc=2pt,
  left=6pt, right=6pt, top=4pt, bottom=4pt, title=#1,
  fonttitle=\bfseries, coltitle=white
}

% Ramka na pułapkę z kartkówki
\newtcolorbox{pulapka}[1][]{
  colback=pulapka, colframe=okres, boxrule=0.8pt, arc=2pt,
  left=6pt, right=6pt, top=4pt, bottom=4pt, title=#1,
  fonttitle=\bfseries, coltitle=white
}

% Nagłówek karty: numer, tytuł, cel
\newcommand{\naglowekKarty}[3]{%
  {\LARGE\bfseries\color{cel} Karta #1: #2}\\[2pt]
  {\large Cel: #3}\par
  \vspace{0.3em}\noindent\rule{\linewidth}{1pt}\par
}

% Pole na obliczenia ucznia: ramka z delikatną kratką, argument = wysokość
\newcommand{\miejsce}[1]{%
  \par\vspace{0.2em}\noindent
  \begin{tikzpicture}
    \draw[step=0.5cm, gray!25, very thin] (0,0) grid (\textwidth,#1);
    \draw[gray!60, rounded corners=2pt] (0,0) rectangle (\textwidth,#1);
    \node[anchor=north west, font=\scriptsize, text=gray] at (0.05,#1) {obliczenia};
  \end{tikzpicture}\par}
 \begin{document}

% Silnik tectonic (XeTeX): polskie znaki przez fontspec, nie przez inputenc/T1
\usepackage{fontspec}
\setmainfont{Arial}
\usepackage[polish]{babel}
\usepackage[a4paper,margin=1.5cm]{geometry}
\usepackage{amsmath,amssymb}
\usepackage{tikz}
\usetikzlibrary{arrows.meta,positioning,calc,shapes.geometric,decorations.pathreplacing}
\usepackage[most]{tcolorbox}
\usepackage{enumitem}
\usepackage{booktabs}
\usepackage{xcolor}

\pagestyle{empty}
\setlength{\parindent}{0pt}
\setlength{\parskip}{0.4em}

\definecolor{cel}{HTML}{1F4E79}
\definecolor{regula}{HTML}{E8F1FA}
\definecolor{pulapka}{HTML}{FDECEA}
\definecolor{okres}{HTML}{C0392B}
\definecolor{skonczony}{HTML}{1E8449}

% Ramka na regułę
\newtcolorbox{regula}[1][]{
  colback=regula, colframe=cel, boxrule=0.8pt, arc=2pt,
  left=6pt, right=6pt, top=4pt, bottom=4pt, title=#1,
  fonttitle=\bfseries, coltitle=white
}

% Ramka na pułapkę z kartkówki
\newtcolorbox{pulapka}[1][]{
  colback=pulapka, colframe=okres, boxrule=0.8pt, arc=2pt,
  left=6pt, right=6pt, top=4pt, bottom=4pt, title=#1,
  fonttitle=\bfseries, coltitle=white
}

% Nagłówek karty: numer, tytuł, cel
\newcommand{\naglowekKarty}[3]{%
  {\LARGE\bfseries\color{cel} Karta #1: #2}\\[2pt]
  {\large Cel: #3}\par
  \vspace{0.3em}\noindent\rule{\linewidth}{1pt}\par
}

% Pole na obliczenia ucznia: ramka z delikatną kratką, argument = wysokość
\newcommand{\miejsce}[1]{%
  \par\vspace{0.2em}\noindent
  \begin{tikzpicture}
    \draw[step=0.5cm, gray!25, very thin] (0,0) grid (\textwidth,#1);
    \draw[gray!60, rounded corners=2pt] (0,0) rectangle (\textwidth,#1);
    \node[anchor=north west, font=\scriptsize, text=gray] at (0.05,#1) {obliczenia};
  \end{tikzpicture}\par}


\begin{document}

\naglowekKarty{3}{Obniżki i podwyżki: mnożnik}{obliczyć nową cenę, także po dwóch zmianach z rzędu}

\vspace{0.3em}
\begin{regula}[REGUŁA – MNOŻNIK ZMIANY CENY]
Po obniżce o $p\%$ zostaje $(100-p)\%$ ceny, po podwyżce o $p\%$ jest $(100+p)\%$ ceny:
$$\text{obniżka o } 20\% \to \cdot\, 0{,}8 \qquad\qquad \text{podwyżka o } 15\% \to \cdot\, 1{,}15$$
\textbf{Dwie zmiany z rzędu:} mnożniki mnożymy, procentów \emph{nie} dodajemy ani nie odejmujemy.
\end{regula}

\vspace{0.4em}
\textbf{Krok po kroku: bluza za $120\ \text{zł}$, obniżka o $25\%$.}
\begin{enumerate}[leftmargin=*, itemsep=3pt]
  \item Zostaje $75\%$ ceny, czyli mnożnik $0{,}75$.
  \item $120\cdot 0{,}75=90\ \text{zł}$.
\end{enumerate}
\textbf{Druga droga:} policz zmianę klockami i odejmij. $25\%$ z $120=30$, więc $120-30=90\ \text{zł}$. Dla podwyżki dodajemy: bilet $40\ \text{zł}$ po podwyżce o $10\%$ to $40+4=44\ \text{zł}$.

\begin{center}
\begin{tikzpicture}[>=Stealth]
  % Skala: 1 cm = 10 zl
  \fill[cel!30] (0,1.6) rectangle (10,2.3);   \draw[thick] (0,1.6) rectangle (10,2.3);
  \fill[okres!30] (0,0.8) rectangle (8,1.5);  \draw[thick] (0,0.8) rectangle (8,1.5);
  \fill[skonczony!30] (0,0) rectangle (9.6,0.7); \draw[thick] (0,0) rectangle (9.6,0.7);
  \node[font=\small\bfseries, anchor=west] at (10.2,1.95) {$100\ \text{zł}$};
  \node[font=\small\bfseries, anchor=west] at (8.2,1.15) {$80\ \text{zł}$};
  \node[font=\small\bfseries, anchor=west] at (9.8,0.35) {$96\ \text{zł}$};
  \node[font=\small, anchor=east] at (-0.2,1.55) {$-20\%$ $\downarrow$};
  \node[font=\small, anchor=east] at (-0.2,0.75) {$+20\%$ $\downarrow$};
  \node[font=\small\bfseries, color=okres] at (5,-0.5) {Nie wróciliśmy do $100\ \text{zł}$: $\ 1{,}2\cdot 0{,}8=0{,}96$};
\end{tikzpicture}
\end{center}

\vspace{0.2em}
\textbf{Przykłady rozstrzygnięte:}
\begin{center}
\renewcommand{\arraystretch}{1.5}
\begin{tabular}{@{}lll@{}}
\toprule
Zadanie & Obliczenie & Wynik \\
\midrule
$120\ \text{zł}$, obniżka $25\%$ & $120\cdot 0{,}75$ & $90\ \text{zł}$ \\
$40\ \text{zł}$, podwyżka $10\%$ & $40\cdot 1{,}1$ & $44\ \text{zł}$ \\
$200\ \text{zł}$, $-10\%$, potem $+10\%$ & $200\cdot 0{,}9\cdot 1{,}1=180\cdot 1{,}1$ & $198\ \text{zł}$ \\
$100\ \text{zł}$, $+50\%$, potem $-50\%$ & $100\to 150\to 75$ & $75\ \text{zł}$ \\
Po obniżce $20\%$ cena $80\ \text{zł}$ & $80\%$ ceny to $80$, więc $100\%$ to $100$ & $100\ \text{zł}$ \\
\bottomrule
\end{tabular}
\end{center}

\begin{pulapka}[PUŁAPKI Z KARTKÓWKI]
\begin{enumerate}[leftmargin=*, itemsep=2pt]
  \item Obniżka o $10\%$, potem podwyżka o $10\%$ \emph{nie} przywraca ceny: dostajemy $99\%$ ceny, nie $100\%$. Drugi procent liczysz z \emph{nowej} ceny.
  \item Cena po obniżce o $20\%$ wynosi $80\ \text{zł}$. Cena przed to $100\ \text{zł}$, a nie $80+20\%=96\ \text{zł}$.
  \item Pytają „o ile zł” czy „ile kosztuje po”? Obniżka to $30\ \text{zł}$, cena po obniżce to $90\ \text{zł}$.
\end{enumerate}
\end{pulapka}

\end{document}
